\section{Relationship to existing evidence and methods}
\label{sec:related}

Split learning was originally introduced to train collaborative models
without exporting raw client records \cite{vepakomma2018}, and recent
surveys catalogue a rapidly growing literature adapting split architectures
to LLM fine-tuning \cite{splitsurvey2026}. Backward tensors still leak
text. Modern reconstruction attacks recover private inputs from exchanged
gradients and activations \cite{gisurvey2022,dager2024,fedspy2026}. Within
our examined programme, adding returned gradients in a controlled GPT-2
benchmark lifts exact-document recovery from $13.71\%$ to $37.77\%$ while
increasing token recovery by $3.17\pp$ \cite{gradients2026}, and
client-returned loss gradients nullify synthetic decoys through exact-zero
support masks \cite{paper1}. Temporal gradient correlations leak
trajectories through similar structural channels in reinforcement learning
\cite{trace2026}. Here we retain each cited measurement under its original
attribution and denominator.

Methodological overlap with LLM-Attack Suite calls for clear boundaries
\cite{suite2026}. That framework already specifies observation contracts,
undefended controls, saved predictions, missing-attack disclosures, and
harness checks. Our work operates one layer higher. By reconciling
cross-campaign decision tables where qualified breaches, uncalibrated
probes, shifted labels, and host-key assumptions collide, we isolate which
programme-level conclusions survive without claiming diagnostic benchmarks
on unseen external repositories.

External studies warn repeatedly that weak evaluators manufacture false
privacy. Wu et al.\ \cite{adaptive2022} break federated defences calibrated
against static attacks, while Smith et al.\ \cite{mia2026} measure
$1.16$-fold to $6.59$-fold higher leakage on FaceScrub under adaptive
attack optimisation (and separately show that changing the external
evaluator's feature basis alters measured leakage). Threat-model and budget
sensitivity appears across multimodal and LLM benchmarks as well
\cite{labelbudget2026,ablation2026}. Layer-wise analyses independently
confirm that deeper transformer splits retain recoverable internal states
\cite{depth2025,actinv2026,bettisplit2026}. Evaluator scepticism still cuts
only one way, because disqualifying 20 uncalibrated probes in our
\RoneCells-cell exploration leaves all \RoneQualifiedBreaches\ breaches
caught by the four calibrated probes intact.

Recent defence studies reinforce both our host-custody and
variational-bound findings. Jin et al.\ \cite{privpair2026} show that
perturbation alone struggles during federated split-LLM fine-tuning and
propose client obfuscate-and-recover adapters, Chen et al.\
\cite{mixguard2026} introduce mixup-based activation obfuscation for split
LLMs, and Zhang et al.\ \cite{aegis2026} investigate generative gradient
masking in medical federated learning. Key custody remains decisive. Thomas
et al.\ \cite{permattack2025} invert permutation-based third-party LLM
inference schemes and illustrate why invertible coordinate transforms
demand explicit host-memory accounting. On the information-theoretic side,
Noorbakhsh et al.\ \cite{mitech2026} apply mutual-information
regularisation during collaborative split inference and Gu et al.\
\cite{pidi2026} incorporate a mutual-information regulariser during
split-LLM fine-tuning, while formal statistical limits on neural
mutual-information estimators \cite{belghazi2018,mcallester2020,song2020},
paired with empirical differential-privacy auditing \cite{jagielski2020},
explain why near-zero critic scores never certify small upper leakage
budgets.

\section{Limitations and implications for further evaluation}
\label{sec:limits}

\paragraph{Selection and dependence.}
Case selection occurred retrospectively after all campaign outcomes had
already been observed. Shared codebases, corpora, checkpoints, and design
decisions link every episode inside one group. Archive tallies describe
this collection alone, and measuring false-positive or false-negative
classification rates on external papers would require a prospectively
registered validation cohort.

\paragraph{Evidence available for recomputation.}
Executable scripts recompute seed contrasts, confidence bounds, and
composite gate decisions wherever raw CSV or JSON arrays exist in the
package. Some upstream artifacts remain source-reported. Decoy-gradient
frame counts and portions of the boundary-replacement review come from
their primary manuscripts and archived review verdicts without refitting
attackers or rescoring tensor captures. File-level SHA-256 digests verify
byte identity across all \PinnedInputs\ consumed records, and untraced
constants lacking a primary calibration curve stay excluded.

\paragraph{Qualification and uncertainty.}
Probe qualification binds strictly to one declared target, channel, defence
mechanism, and budget. Seed cohorts remain small. Four-seed two-sided
Student-$t$ intervals in the exploration and three-seed one-sided upper
bounds in the defence pilot condition on fixed evaluation panels without
multiplicity adjustments, while the unexecuted eight-seed power figure
serves purely as a variance-sensitive planning bound. Positive-control
canary floors, permutation-null bands, acceptance gates, and minimum
detectable leak rates measure distinct properties even when two constants
share the value $0.02$.

\paragraph{Security and utility.}
Empirical non-detection reflects a bounded attacker and compute budget. Low
pooled passive recovery, invertible wrappers whose keys sit in cloud RAM,
and near-zero variational lower bounds leave compromised-host secrecy
unproved. Utility gates cover only their scored endpoints, since
short-context cross-entropy can pass while teacher-forced answer-token
accuracy drops sharply, and evaluating noise or quantisation fairly
requires sweeping perturbation strength against explicit task-quality caps.

Explicit threat contracts enable testable future benchmarks. Deliberately
injected channel omissions and label shifts can run alongside genuinely
leaky defences under frozen compute budgets. Prospective validation trials
on held-out model families and corpora can then quantify how reliably
automated admissibility checks catch false privacy assurance while
preserving confirmed breaches.

\section{Conclusion}

Re-examining six linked investigations from one split-LLM programme
separates broken measurements from broken defences. Two early assurances
failed on method. Omitted client-returned gradients and off-by-one token
labels undermined those initial claims, yet \RoneQualifiedBreaches\
qualified breaches in the \RoneCells-cell exploration alongside $40\times$
punctuation-property excesses and $5\times$ next-byte structure excesses in
the seven-cell pilot stand as sound negative evidence against the tested
defences despite incomplete probe qualification elsewhere. Scope checks
likewise confine shared-key orthogonal wrappers to passive network links,
document the $63.16\%$ throughput penalty of rapid per-epoch ratcheting,
and show why variational mutual-information lower bounds near zero leave
true leakage unconstrained above. Tying every privacy verdict to its
observed channels, key custody, calibrated controls, and seed-level
uncertainty establishes exactly what archived split-training experiments
support.

{\emergencystretch=2em
\paragraph{Data and code availability.}
The manuscript source, reanalysis scripts, derived tables, and archived
evidence reside in the companion repository at \repourl{}
\cite{artifacts2026}. Python recomputation scripts live under
\path{analysis/}, derived tables and SHA-256 input manifests under
\path{data/}, LaTeX sources and vector figure generators under
\path{paper/}, and primary records under \path{evidence/} and
\path{reproduction/}. Appendix~\ref{app:artifacts} maps each retained
statistic to its exact repository-relative path. Every consistency check
runs offline on the Python standard library without GPU retraining.\par}
